\section{Q\&A: Restaurar el eslogan de arquitectura en juicios de ingeniería}

En la sesión de Q\&A, aproximadamente 23 minutos después del discurso principal, no se añadieron diapositivas; la cámara continuó enfocando al ponente y al fondo con conclusiones. Sin embargo, contenía las condiciones de borde más importantes de esta lección: cómo mantener la fusión mediante especialización, cómo auditar por separado los objetivos y los parámetros, qué ideas solo merecen ser experimentadas y si el lenguaje es necesariamente un medio para el razonamiento. Este capítulo reorganiza estas respuestas bajo preguntas de ingeniería, en lugar de copiarlas pregunta por pregunta.

\subsection{¿Cómo se fusiona la información una vez separados los parámetros?}

La proyección específica de modalidad (modality-specific projection) de MoT no canceló la atención de secuencia mixta (mixed-sequence attention). Si definimos la consulta de texto como $q_T=W_Q^T h_T$ y la clave de imagen como $k_V=W_K^V h_V$, aún podemos calcular:
$$
\operatorname{score}(T,V)=\frac{q_T^\top k_V}{\sqrt{d}}.
$$
Aunque $W_Q^T$ y $W_K^V$ pertenecen a familias de parámetros diferentes, el producto punto sigue integrando ambas modalidades en un mismo grafo de atención. Por ello, el modelo puede «codificar de manera especializada y leer de manera conjunta».

\teachervoice{00:42:42--00:44:08, el profesor respondió a la pregunta «¿bloquea el aislamiento de parámetros el intercambio de información?» utilizando la atención por sí misma: los tokens pasan por parámetros distintos, pero eso no impide que se vean mutuamente en la atención. Al leer un diagrama de arquitectura, hay que rastrear por separado la ruta de parámetros y la ruta de información.}

\begin{table}[H]
\centering
\small
\begin{tabularx}{0.98\textwidth}{L{0.23\textwidth}Y Y}
\toprule
Eje de diseño & Opciones & Pregunta central \\
\midrule
Representación & Códigos VQ, embedding semántico, latente VAE, codec de audio & ¿Qué información se conserva en la entrada/salida? \\
Objetivo & CE, difusión, coincidencia de flujo, pérdida de acción & ¿Qué predice el modelo para ser recompensado? \\
Compartición de parámetros & totalmente compartido, específico de modalidad, adaptador, MoE & ¿Qué tokens compiten por la misma capacidad? \\
Mezcla de información & atención global, atención cruzada, fusión tardía & ¿En qué capa intercambian las modalidades la información? \\
Servicio & síncrono, transmisión en tiempo real, razonamiento en segundo plano & ¿Cómo se organiza la latencia y la interacción con el usuario? \\
\bottomrule
\end{tabularx}
\caption{Ejes de diseño que deben auditarse por separado en sistemas multimodales}
\end{table}

\subsection{La elección del objetivo y la compartición de parámetros no son lo mismo}

Un backbone compartido puede optimizar simultáneamente la entropía cruzada (CE) y la difusión; un conjunto de parámetros específico de modalidad también puede generar texto. El valor del Q\&A radica en rechazar la dicotomía de «ver compartido y asumir el mismo objetivo» o «ver diferentes objetivos y asumir dos modelos». La implementación de ingeniería debe registrar para cada lote (batch) las máscaras de modalidad, la normalización de la pérdida, la norma del gradiente y los parámetros activados, solo así se puede ubicar de dónde proviene la interferencia.

\begin{importantbox}{Superficie mínima de observación para el entrenamiento multitarea}
Registrar por cada modalidad: cantidad de muestras, cantidad de tokens, pérdida, norma del gradiente, parámetros activos, rendimiento y métrica de validación; además, registrar sondas de atención/representación entre modalidades. Sin esta descomposición, la disminución del objetivo global de entrenamiento podría ocultar el continuo deterioro de una modalidad específica.
\end{importantbox}

\subsection{¿La generación de video es un modelo del mundo?}

El Q\&A discutió si generar fotogramas futuros puede ayudar a la predicción de acciones. Si el modelo puede predecir $o_{t+1}$ bajo una acción condicional, efectivamente ha aprendido parte de las dinámicas del entorno; sin embargo, un «video que parece razonable» difiere aún de un «modelo de transición calibrado para control». Este último debe mantener consistencia ante acciones contrafácticas, fallos raros y despliegues a largo plazo (long-horizon rollout).

\teachervoice{00:46:16--00:48:00, el profesor consideró la generación de video como un posible retroalimentación predictiva, en lugar de declarar directamente que los modelos generativos son modelos del mundo confiables. Las notas conservan este matiz condicional: se requiere retroalimentación de tareas y verificación en bucle cerrado.}

\begin{warningbox}{Tres comprobaciones para la afirmación de modelo del mundo}
1. Condicionamiento por acción: ¿la predicción es igual si no cambiamos la acción? 2. Consistencia del despliegue (rollout): ¿el error explota con el tiempo? 3. Utilidad de decisión: ¿el uso de las predicciones mejora realmente el éxito de la política? Si solo hay fotogramas futuros bonitos sin estos tres puntos, como máximo demuestra capacidad de generación de video.
\end{warningbox}

\subsection{«Dibujar texto en imagen» es un problema experimental, no una conclusión establecida}

Un estudiante preguntó si era posible renderizar texto como un mapa de bits (bitmap) blanco y negro para modelar todo mediante una representación visual unificada. El ponente aclaró que no lo había hecho personalmente y juzgó que los tokens simbólicos limpios podrían ser más eficientes. Este tipo de respuesta vale la pena conservarlo porque ejemplifica el correcto enfoque hacia ideas novedosas: definir primero experimentos de control, en lugar de empaquetar intuiciones como hechos.

Un experimento razonable debe fijar el corpus, las FLOPs del modelo y la información de contexto para comparar: tokens de texto originales; parches de mapa de bits; codificador de mapa de bits preentrenado con OCR; y tokens híbridos. Además de la perplejidad, se debe evaluar la longitud de secuencia, el rendimiento de entrenamiento, la transferencia de razonamiento y la robustez. El mapa de bits podría aumentar el compartimiento visual, pero también introducir redundancia tipográfica, de renderizado y espacial.

\teachervoice{00:48:14--00:53:31, el profesor marcó repetidamente «nunca lo he probado personalmente» y «vale la pena intentar experimentos», al final solo ofreciendo intuición sobre eficiencia. Las notas no actualizan esta respuesta como una ley de escala (scaling law).}

\subsection{Representación centrada en objetos y jerarquía}

La cuadrícula de parches facilita los cálculos, pero no implica que el mundo esté compuesto por bloques fijos. El Q\&A propuso un embedding orientado a objetos, una representación jerárquica y unidades de escena más semánticas. Intentan permitir que el modelo represente «objetos—relaciones—eventos», reduciendo la sensibilidad a cambios a nivel de píxel. El desafío abierto es descubrir objetos estables en datos no supervisados, manejar oclusiones y continuidad de identidad, mientras se apoya la generación.

\begin{knowledgebox}{Estructura en tres niveles del frente de representación}
El nivel bajo conserva textura, frecuencias sonoras y movimiento; el medio forma objetos, partes, oradores y eventos; el alto forma relaciones, intenciones y estado de tarea. Las tareas de comprensión prefieren invarianzas en el nivel alto, mientras que las de generación requieren detalle en el nivel bajo. Un sistema escalable probablemente use una jerarquía, en lugar de resolver todas las tareas en una única resolución.
\end{knowledgebox}

\subsection{El lenguaje es la armazón actual del razonamiento, pero no un requisito lógico}

La última pregunta extensa señaló que los datos sintéticos de texto son fáciles de generar y pueden ocultar el pensamiento en cadena (chain-of-thought), mientras que los fotogramas de razonamiento visual son costosos y el usuario los vería directamente. El ponente estuvo de acuerdo en que el lenguaje es actualmente una armazón de abstracción muy efectiva, y también aceptó que la generación visual está limitada por eficiencia y experiencia de usuario (UX); sin embargo, no descartó el razonamiento en un espacio puramente visual en el futuro. Si un modelo predice con precisión futuros complejos de video y apoya la toma de decisiones, también puede considerarse una forma de razonamiento no textual.

\teachervoice{01:03:32--01:04:21, la postura final del profesor es de «doble limitación»: la armazón lingüística es muy útil ahora; el razonamiento puramente visual sigue siendo una pregunta abierta en principio. No se debe reescribir las ventajas actuales de ingeniería como leyes de teoría cognitiva.}

\begin{warningbox}{Límites profesionales para respuestas sobre razonamiento espacial}
00:58:34, el profesor aclaró explícitamente que no es experto en razonamiento espacial, sino que solo generalizó los avances recientes en robótica. Por ello, las notas no presentan esta respuesta como un repaso del estado de la técnica (state-of-the-art survey); solo apoya un juicio a nivel de curso: existe una distancia significativa entre las capacidades multimodales numéricas y el razonamiento espacial confiable.
\end{warningbox}

\subsection{Resumen de la sección}

El Q\&A descompone «unificado» y «nativo» en variables de ingeniería: los parámetros pueden ser especializados mientras la información se mantiene fusionada; el objetivo y el backbone son ejes independientes; la generación de video se acerca al modelo del mundo solo si es efectiva en la toma de decisiones; las hipótesis de nuevas representaciones requieren experimentos pareados; el lenguaje es actualmente una interfaz de razonamiento eficiente, pero sus formas futuras siguen siendo abiertas.
